\chapter{微分几何：弯曲空间中的距离和方向}

\section{地图上那条奇怪的弯路}

打开任意一张世界地图，找到北京和旧金山。现在想象你是一名民航飞行员，要规划两地之间最省油的航线。直觉告诉你：拿一把直尺，在地图上把两个城市连起来，沿直线飞就是了。

可是真实的航班不这么飞。从北京起飞的航班会先向东北方向爬升，掠过我国东北和俄罗斯远东，擦过阿留申群岛附近，再折向东南降落在旧金山。在地图上看，这条路线明显是弯的——它向北鼓出一个大大的弧线。航空公司每年烧掉的航空燃油以万吨计，绝不可能无缘无故绕远路。

再来看一个更古老的例子。十五世纪的航海家早就发现：沿着同一条纬线一直向西航行，并不是去远方港口最省时间的路；聪明的船长会先向高纬度方向偏航，再转回来。船员们管这叫"大圆航行"。他们不知道为什么，只知道这样走更快。

于是问题摆在了我们面前：

\begin{center}
\begin{tikzpicture}
  \draw[thick] (0,0) circle (2.2);
  \coordinate (P) at (0,2.2);
  \coordinate (A) at (-2.0,-0.9);
  \coordinate (B) at (2.0,-0.9);
  \draw[thick] (A) to[out=70,in=180] (P);
  \draw[thick] (P) to[out=0,in=110] (B);
  \draw[thick] (A) to[out=-15,in=195] (B);
  \fill (P) circle (0.06) node[above] {$A$};
  \fill (A) circle (0.06) node[left] {$B$};
  \fill (B) circle (0.06) node[right] {$C$};
  \node at (0,-3.2) {球面上的一个三角形：三条边都是大圆弧};
\end{tikzpicture}
\end{center}

在弯曲的表面上，"直"是什么意思？"距离"该怎么量？"方向"又该怎么描述？这些问题听起来像是给飞行员和船长准备的，但它们背后是一门深刻的数学——微分几何。它研究弯曲空间中的长度、角度和形状，并且在二十世纪初成为爱因斯坦描述引力的语言。这一章，我们从一条绕弯的航线出发，一直走到弯曲的时空。

\section{直觉失灵：摊不平的橘子皮}

面对球面上的距离问题，第一个自然想法是：把球面摊平不就行了？摊成平面，用尺子量，多省事。

你可以在家做个实验。剥一个橘子，尽量保持橘皮完整，然后试着把橘皮平平地铺在桌面上，不许有裂缝也不许重叠。你会发现根本做不到：要么边缘裂开一道道口子，要么中间拱起来。现在对比另一件事：拿一张平整的纸，把它卷成一个圆筒，毫不费力；展开，纸又恢复原样，没有撕裂也没有褶皱。

这个厨房里的失败实验其实揭示了一条深刻的分界线：圆柱面和平面是"亲戚"，可以不拉伸地互相变形；球面和平面不是。任何把球面画到平面上的地图，都必然要付出扭曲的代价。

代价有多大？看看最常见的世界地图（墨卡托投影）：地图上格陵兰岛看起来和非洲差不多大，但实际上非洲的面积大约是格陵兰的十四倍。地图没有"画错"，它只是不可能同时保住所有的长度、角度和面积。有些地图保住角度就牺牲面积，有些保住面积就牺牲形状。数学家们证明了：一张既保距又保角又保面积的球面地图，根本不存在。

\begin{fanli}
"把球面摊平了再量，不就回到平面几何了吗？"——这是一个看似合理、实则走不通的方案。球面无法在不拉伸、不撕裂的情况下摊成平面，这意味着球面上的几何学不是平面几何的"化装舞会"，而是一门有自己规律的独立学问。同理，"沿纬线飞就是从北京到旧金山的最短路"也是错的：除了赤道，任何纬线圈都不是大圆，沿纬线走一定绕远。
\end{fanli}

第一次尝试失败了，但失败指了路：我们不能再想着"把曲面变成平面"，而要学会在曲面内部生活——像球面上的一只蚂蚁那样，不跳出来，只靠脚下的测量理解整个世界。

\section{给曲面建立坐标：一本地图册就够了}

蚂蚁不需要一张覆盖整个球面的完美地图。它需要的只是一本地图册：每一页只画球面上的一小片区域，这一小片可以近似当成平面；相邻的两页之间有一段重叠，重叠处的同一个地点在两页上有两套坐标，并且附有一份"换算规则"，告诉你怎么把一套坐标翻译成另一套。

这个想法非常重要，值得给它一个名字。

\begin{dingyi}
曲面上的一小块区域，如果能和平面上的一个圆盘一一对应、并且对应关系足够平滑，我们就给它配上平面坐标 $(u,v)$，称为一张\textbf{坐标图}。若干张坐标图拼起来、彼此重叠处可以互相换算，覆盖整个曲面，合起来称为一册\textbf{坐标图册}。
\end{dingyi}

你在家里其实就有这样的东西。没有哪张世界地图能完美呈现整个地球，但一本分省地图册很好用：每页一个省，页边的重叠区域帮你翻页时接上。球面也是这样：一张坐标图覆盖不了全球（极点处总会出麻烦），但两三张图拼起来就够了。关键是，几何量——长度、角度、面积——必须在每张图里都能算，而且不同图算出来的结果在重叠处必须一致。只要满足这一点，蚂蚁就拥有了完整的几何学，它根本不需要知道"球面嵌在三维空间里"这回事。

这种"只在曲面内部看问题"的视角，叫作内蕴视角。与之相对的是站在外面看：三维空间里的球面当然是弯的，一眼就能看出来。数学史上的一个转折点，是数学家们意识到内蕴几何可以独立存在：曲面自己弯不弯、弯多少，不用到"外面"去也能测量出来。一位数学家把这方面的核心结果称为"绝妙定理"——内蕴曲率是曲面的固有属性，压扁、卷起这些不改变内蕴曲率的操作，骗不过生活在曲面上的测量者。

\section{切平面与切向量：在曲面上站稳}

蚂蚁要在球面上爬行，首先得回答一个问题：什么叫"方向"？

在平面上，方向就是箭头——向量。可球面上的箭头往哪儿指？如果箭头要画在球面上，它得弯曲地贴着球面，那就不是我们熟悉的向量了。解决办法是：在球面上的点 $P$ 处，放一块无限薄的平板，让它刚好只在 $P$ 点接触球面，其他地方都"悬空"。这块平板叫作球面在 $P$ 点的\textbf{切平面}。

\begin{center}
\begin{tikzpicture}
  \draw[thick] (0,0) circle (2);
  \draw[thick] (-2.8,2) -- (2.8,2);
  \fill (0,2) circle (0.06) node[above right] {$P$};
  \draw[->,thick] (0,2) -- (1.6,2) node[right] {$v$};
  \node at (0,-2.8) {圆在 $P$ 点的切线：曲面在 $P$ 点的切平面的二维类比};
\end{tikzpicture}
\end{center}

\begin{dingyi}
曲面在点 $P$ 的\textbf{切平面}，是在 $P$ 点与曲面相切、并且包含曲面在该点所有方向的平面。切平面中以 $P$ 为起点的向量，称为曲面在 $P$ 点的\textbf{切向量}。
\end{dingyi}

切向量解决了一个真实的麻烦：它让"在曲面上朝某个方向走一小步"有了精确含义。走的那一小步确实贴着曲面弯曲，但步子无穷小的时候，它和切平面上的直线段几乎没有差别——误差比步长本身缩小得更快。这正是"微分"两个字的分工：无穷小的局部，弯曲的表面和它的切平面是同一种东西。

你其实每天都活在这个近似里。地球半径约 $6371$ 千米，你脚下的地面是球面的一部分；可是在方圆几十米的尺度上，大地和切平面的偏差只有毫米量级，盖房子、踢足球完全可以当它是平的。局部像平面，整体是弯的——请记住这句话，它后面还会以更重要的身份出现。

\begin{shuyu}{切平面}
\textbf{直观解释}：把一块玻璃板轻轻贴在篮球表面的某一点上，玻璃板所在的平面就是切平面。换一点，玻璃板的朝向就变了。

\textbf{正式定义}：曲面在一点的切平面，是该点处所有光滑曲线的切线（方向）张成的平面。

\textbf{一个例子}：地球表面上，你所站的地面附近那一小块，就近似是你脚下那点的切平面；水平仪测的就是它。
\end{shuyu}

\section{曲率：弯了多少，用什么量}

有了切平面，就可以量化"弯"了。想法很直接：曲面上的点离开切平面的速度越快，曲面在那里就越弯。

先看一维的类比。平面上半径为 $r$ 的圆，它的弯曲程度处处相同，数学家定义它的曲率为 $1/r$：半径越小，圆弯得越急；半径趋于无穷大时，曲率趋于零，圆就退化成了直线——直线的曲率是零。数值例子：半径 $r=2$ 的圆，曲率是 $1/2$；半径 $r=100$ 的圆，曲率只有 $0.01$，你在上面走一小段几乎感觉不出它在弯。

曲面上的情况复杂一点：在同一点，朝不同方向走，弯曲的程度可能不一样。球面上每个方向都一样弯，半径为 $R$ 的球面，沿任何方向的截线都是半径 $R$ 的大圆，各个方向的曲率都是 $1/R$。圆柱面就不同了：沿着圆环方向截，得到一个半径 $R$ 的圆，曲率 $1/R$；沿着母线（竖直方向）截，得到一条直线，曲率是零。把两个极端方向的曲率乘起来，得到的量叫\textbf{高斯曲率}。于是：

\[ \text{球面的高斯曲率} = \frac{1}{R}\cdot\frac{1}{R} = \frac{1}{R^2}, \qquad \text{圆柱面的高斯曲率} = \frac{1}{R}\cdot 0 = 0. \]

这个结果初看违反直觉：圆柱面明明"看起来是弯的"，高斯曲率怎么是零？回想剥橘皮的实验就明白了：圆柱面可以由一张纸不拉伸地卷成，生活在圆柱面上的蚂蚁做的任何测量，和平面上的蚂蚁完全一样——它的内蕴几何就是平面几何。高斯曲率量的正是这种内蕴的弯，而不是从外面看起来的弯。反过来，球面的高斯曲率 $1/R^2$ 处处为正，这就是橘皮摊不平的定量证据。数值例子：地球半径约 $6371$ 千米，地表的高斯曲率约为 $1/(6.371\times 10^6)^2 \approx 2.5\times 10^{-14}$ 每平方米——极其微小的正数，但不为零，日积月累就让你的世界地图永远画不完美。

还有第三种可能：马鞍面。在马鞍的中心点，沿前后方向曲面向下弯，沿左右方向向上弯，两个方向的弯曲"符号相反"，乘出来的高斯曲率是负的。负曲率的曲面摊到平面上会怎样？它既不像球面那样"料不够"（会裂），而是"料太多"——摊平后会起皱、重叠，像一片木耳。三种符号，三种命运：正曲率裂开，零曲率服帖，负曲率起皱。

\section{球面三角形：内角和的反叛}

初中几何的第一批定理里，有一条刻在我们脑子里：三角形内角和等于 $180^\circ$。现在我们要看着它在球面上失效，并且失效得很有规律。

先明确球面上"三角形"的定义。球面上扮演直线角色的是大圆——过球心的平面与球面的交线，比如赤道和所有的经线圈（纬线圈除了赤道都不是大圆，因为它们所在的平面不过球心）。

\begin{dingyi}
球面上取三个不在同一大圆上的点，用三段大圆弧把它们两两连接，围成的图形叫作\textbf{球面三角形}。它的三个内角，是顶点处两段大圆弧的切线之间的夹角。
\end{dingyi}

立刻看一个极端又漂亮的例子。取北极点 $A$，再取赤道上经度相差 $90^\circ$ 的两个点 $B$、$C$。边 $AB$、$AC$ 都是经线，边 $BC$ 是赤道的一段。顶点 $A$ 处，两条经线的夹角就是经度差 $90^\circ$；顶点 $B$ 处，经线与赤道垂直，夹角 $90^\circ$；顶点 $C$ 同理，也是 $90^\circ$。三个角都是直角，内角和是 $270^\circ$，比平面上整整多出 $90^\circ$！

这多出来的部分不是误差，而是面积。十九世纪初，一位数学家证明了球面几何里最优美的公式之一：

\begin{dingli}[球面三角形的面积公式]
设球面半径为 $R$，球面三角形的三个内角（用弧度制）为 $A$、$B$、$C$，面积为 $S$，则
\[ S = R^2\,(A + B + C - \pi). \]
特别地，球面三角形的内角和一定大于 $\pi$（即 $180^\circ$），且超出的量 $A+B+C-\pi$ 与面积成正比。
\end{dingli}

\begin{proof}
证明分三步，每一步都可以在图上指出来。

\textbf{第一步：二角形的面积。} 两个大圆相交，把球面分成四个"月牙"形的区域，每个叫作二角形。夹角为 $A$（弧度）的二角形，占整个球面的比例是 $A/(2\pi)$（把球面像切西瓜一样按经度均分成 $2\pi$ 弧度，它占其中 $A$ 弧度）。全球面积为 $4\pi R^2$，所以夹角 $A$ 的二角形面积为
\[ 4\pi R^2 \cdot \frac{A}{2\pi} = 2R^2 A. \]
检验：当 $A=\pi$ 时二角形恰好是半球，公式给出 $2\pi R^2$，与半球面积一致。

\textbf{第二步：三对二角形覆盖整个球面。} 设球面三角形 $T$ 的三个内角为 $A$、$B$、$C$。过 $A$ 的两条边所在的大圆，围出夹角 $A$ 的一对对顶二角形，每个面积 $2R^2 A$；同理得夹角 $B$、$C$ 的各一对二角形。这六个二角形合起来覆盖整个球面，但覆盖有重复：三角形 $T$ 本身被覆盖了 $3$ 次，球面另一侧与 $T$ 对顶的全等三角形也被覆盖了 $3$ 次，其余每点恰好被覆盖 $1$ 次。

\textbf{第三步：列方程。} 六个二角形面积之和，等于全球面积加上多算的 $4$ 份 $T$（$T$ 多算了 $2$ 份，对顶三角形也多算了 $2$ 份）：
\[ 2\,(2R^2 A + 2R^2 B + 2R^2 C) = 4\pi R^2 + 4S. \]
两边除以 $4R^2$，整理得
\[ S = R^2 (A + B + C - \pi). \]
证毕。
\end{proof}

用前面的数值例子验证：北极加赤道两点的三角形，$A=B=C=\pi/2$，内角和为 $3\pi/2$，公式给出 $S = R^2(3\pi/2 - \pi) = \pi R^2 /2$。直接看：这个三角形占全球八分之一（两条相差 $90^\circ$ 的经线切出四分之一，再取北半球的一半），面积 $\frac{1}{8}\cdot 4\pi R^2 = \pi R^2/2$，完全吻合。

\begin{tuilun}
球面上不存在内角和恰好等于 $180^\circ$ 的三角形；内角和超出 $180^\circ$ 越多，三角形越大。当三角形缩得很小时，超出部分趋于零——小范围内，球面几何近似回到平面几何。
\end{tuilun}

这条推论把我们前面的话焊死了：局部像平面，整体是弯的。曲率正是通过"内角和背叛了多少"这种方式，被曲面内部的居民测量出来。

\begin{lianjie}
回忆平面几何中"三角形内角和等于 $180^\circ$"的证明：它依赖平行公设——过直线外一点恰有一条平行线。球面上任意两个大圆都相交，"平行线"根本不存在，平行公设失效，内角和定理随之改写。这不是哪个定理的证明出了错，而是几何的舞台换了。历史上，正是对平行公设两千年的怀疑与争论，最终把数学家们引向了弯曲空间的几何；本章的球面几何，就是那场内战留下的一份遗产。到了下一章，你会看到另一个"习以为常的确定性"——一次测量的精确数字——同样会在随机性面前改写。
\end{lianjie}

\section{测地线：最直的，也是最短的}

现在回到开篇的问题：什么是在曲面上"直走"？蚂蚁没有罗盘，它对"直走"的理解是：每一步都不向左拐也不向右拐，头朝哪边就往哪边爬。这样的路线叫\textbf{测地线}。

\begin{shuyu}{测地线}
\textbf{直观解释}：曲面上"不拐弯"的路线；在曲面上两点间拉紧一根橡皮筋，它静止时贴出的就是测地线。

\textbf{正式描述}：测地线是曲面上局部最短的曲线——它上面充分接近的任意两点之间，没有比它更短的曲面路线；等价地，它是沿曲面"直行"、左右加速度为零的曲线。

\textbf{一个例子}：球面上的测地线就是大圆；北京到旧金山的最短航线，是过两地和地心的大圆上较短的那段弧。
\end{shuyu}

为什么"最直"和"最短"是同一件事？可以这样想：如果一条路线中间有弯，你在弯的地方抄近道切过去，路线就变短了；所以最短的路必然处处不弯。反过来，处处不弯的路，任何一小段都是最短。在平面上，测地线是直线；在球面上，测地线是大圆；在圆柱面上，把圆筒剪开摊平，测地线变成平面上的直线段，卷回去它可能是一条绕筒上升的螺旋线——螺旋也可以是"最直的路线"，这一点第一次听说时总让人愣一下。

来算一笔账，看看大圆航线省多少路。北京约为北纬 $40^\circ$、东经 $116^\circ$，纽约约为北纬 $41^\circ$、西经 $74^\circ$，两地几乎在同一纬度，经度相差约 $170^\circ$。如果"沿纬线直飞"（这正是平面地图上接近直线的走法），走的是北纬 $40^\circ$ 纬线圈，其半径是 $R\cos 40^\circ$，弧长为
\[ \frac{170}{360}\times 2\pi R \cos 40^\circ \approx 0.472 \times 2\pi \times 6371 \times 0.766 \approx 14500 \text{ 千米}. \]
而实际的大圆航线长约 $11000$ 千米，足足省下约三千五百千米——相当于多飞一个北京到乌鲁木齐的来回还有富余。地图上弯的那条弧，才是球面上真正的"直线"。

\begin{shiyan}
找一个橙子或乒乓球、一支记号笔、一根软绳和一把量角器（没有量角器就估角度）。

第一步：用笔在"球"上点三个点，把软绳在两点间拉紧贴住球面，沿绳子画出三段"直线"——你画出的就是大圆弧，即测地线。

第二步：观察这个球面三角形的三个角，估计或量出每个角的度数并相加。把三角形画得越大，内角和超出 $180^\circ$ 越多；画得很小时，内角和接近 $180^\circ$。

第三步（对比实验）：把一张纸上画好的三角形卷成圆筒，再展开量角——内角和纹丝不动。想一想：为什么球面上的三角形"变了"，纸上的却"没变"？
\end{shiyan}

\section{流形：每一小块都是平面，整体可以千奇百怪}

我们已经反复使用同一个句式：局部像平面，整体可能很不一样。数学家把这个句式本身变成了研究对象。

\begin{dingyi}
一个空间，如果它每一点周围的一小块区域都能和平面上的圆盘建立平滑的一一对应（可以画成一张坐标图），就称它为一个\textbf{二维流形}。把"平面"换成 $n$ 维的平直空间，就得到 $n$ 维流形。
\end{dingyi}

\begin{shuyu}{流形}
\textbf{直观解释}：流形就是"地图册能覆盖的空间"。身处其中任何一点，环顾四周你都以为自己站在平地上；只有环球旅行回来，才发现世界另有形状。

\textbf{正式定义}：每一点都有一个邻域与平直空间的开圆盘平滑等价，且图册中地图之间的换算足够光滑。

\textbf{一个例子}：球面是二维流形；甜甜圈表面（环面）也是；而两个平面交叉成的"十字"不是——交叉点的那一小块怎么也摊不成一个圆盘。
\end{shuyu}

环面值得多说两句。甜甜圈表面有的地方像球面（外侧，正曲率），有的地方像马鞍（内侧，负曲率），整体却是一个闭合的、没有边界的流形。生活在环面上的蚂蚁画三角形，会发现有的三角形内角和大于 $180^\circ$，有的小于 $180^\circ$——曲率的符号随地点而变。更有意思的是，把长方形纸的两条对边分别粘合（先卷成筒，再把筒的两头接上），得到的也是一个环面，而这个"纸环面"在内蕴几何上居然是平的——上面居然存在内角和恰为 $180^\circ$ 的几何学。同一个拓扑形状，可以承载不同的几何；这区分"形状"与"几何"两层结构，正是现代几何学的起点。

\begin{table}[h]
\centering
\begin{tabular}{llll}
\toprule
曲面 & 三角形内角和 & 测地线的样子 & 能否不拉伸地摊平 \\
\midrule
平面 & 恰为 $180^\circ$ & 直线 & 本来就是平的 \\
球面 & 大于 $180^\circ$，越大超得越多 & 大圆弧 & 不能（会裂） \\
圆柱面 & 恰为 $180^\circ$ & 展开后是直线（可为螺旋） & 能 \\
马鞍面 & 小于 $180^\circ$ & 向外发散的曲线 & 不能（会皱） \\
\bottomrule
\end{tabular}
\end{table}

流形这个概念为什么值得如此大费周章？因为我们自己就生活在一个"流形"里。抬头看宇宙：我们所能直接经验的，永远是身边一小块近似平直的空间；宇宙整体是闭合的还是开放的、是平的还是弯的，无法一眼看穿，只能靠内部的测量——三角形的内角和、平行光束的聚散、遥远星系的计数——一点点推断。蚂蚁研究球面的方法，正是天文学家研究宇宙的方法。

\section{从局部几何到星空：一次科普性的远眺}

1854 年，一位数学家在演讲中系统提出了"任意维数的弯曲空间"的数学：空间可以在每一点有自己的曲率，长度和角度由一张逐点变化的"度量表"决定，直线让位给测地线。这套工具在当时纯粹是智力游戏——没有人知道它能用来干什么。

六十一年后，1915 年，爱因斯坦发现这正是他需要的语言。广义相对论的核心图景可以这样科普地讲：我们以为的"引力"，其实是时空本身的弯曲；物质和能量的存在让时空弯曲，而弯曲时空中自由下落的物体，沿着时空的测地线运动——行星绕太阳转，并不是被一根无形的绳子拽着，而是在弯曲时空里"直走"。1919 年日全食期间，天文学家观测到遥远星光掠过太阳附近时的偏折，偏折量与理论预言吻合：光线走的是时空的测地线，而太阳让那片时空弯曲了。

必须诚实地标明边界：以上是一段科普类比。真实的广义相对论处理的是四维时空（三维空间加一维时间）的弯曲，其几何与时间的流逝交织在一起，由一组精确的方程描述，远不是"橡皮膜上放铁球"的图片所能概括。本章能给你的，是这幅图景的入口：弯曲的空间可以用内蕴的度量来研究，测地线取代了直线，曲率取代了引力——这三件事，你现在已经能在球面上亲手验证。

\begin{pandeng}
继续攀登的路标：高斯曲率与"绝妙定理"；协变导数与平行移动（向量沿曲线搬运时如何变化）；高斯–博内定理（曲面的总曲率与它的"洞数"挂钩，是几何与拓扑的桥）；黎曼几何与度量张量；再往前，是微分几何与广义相对论的正式课程。这些关键词，等你学了多元微积分和线性代数之后，都可以一一兑现。
\end{pandeng}

\begin{toolbox}
本章新添的工具：

\begin{itemize}
  \item \textbf{坐标图与图册}：用一片片平面坐标覆盖曲面，重叠处负责换算——几何学研究曲面不必把它摊平。
  \item \textbf{切平面与切向量}：曲面在一点的所有方向构成切平面；无穷小尺度下，曲面与切平面无法区分。
  \item \textbf{高斯曲率}：两个极端方向曲率的乘积，刻画内蕴的弯曲；正者摊平会裂，零者服帖，负者起皱。
  \item \textbf{球面三角形面积公式} $S=R^2(A+B+C-\pi)$：内角和的"盈余"就是面积，小三角形近似平面。
  \item \textbf{测地线}：曲面上不拐弯的路，也就是局部最短的路；球面上是大圆。
  \item \textbf{流形}：局部处处像平直空间、整体形状自由的空间——我们自己的宇宙很可能就是一个。
\end{itemize}
\end{toolbox}

\section*{思考题}

\textbf{观察题}

1. 在地球仪（或橙子）上找三个点：北极点，以及赤道上经度相差 $120^\circ$ 的两个点。连出球面三角形，观察它的三个角并相加。

提示：三个角彼此相等；先想顶点处经线之间、经线与赤道之间的夹角各是多少，再对照面积公式检验你的结果。

2. 找一只圆筒形的杯子，在侧面画一个三角形，再把纸面（或想象纸面）摊开。

提示：观察哪些量在卷与摊的过程中保持不变——长度、角度、内角和；想想这说明圆筒面的内蕴几何是什么。

\textbf{证明题}

1. 证明：球面上不存在"内角和等于 $180^\circ$"的退化情形之外的三角形。

提示：直接用面积公式 $S=R^2(A+B+C-\pi)$，注意面积 $S$ 是正数。

2. 证明：球面上任意两个大圆必相交，而且恰好交于一对对径点（直径的两端）。

提示：大圆是"过球心的平面"与球面的交线；两个过球心的平面必交于一条过球心的直线。

\textbf{挑战题}

1. 在球面上，是否存在"三边两两全等、内角和却不同"的两个三角形？在平面上，边边边相等保证全等；球面上还成立吗？

提示：用面积公式把内角和与面积挂钩，再想想球面上三角形的面积有没有上限。

2. 甜甜圈（环面）表面上，有的地方内角和大于 $180^\circ$，有的小于 $180^\circ$。但存在一种"完全平坦"的环面几何（由长方形粘合对边得到）。思考：平坦环面上，测地线有哪些不同的"绕法"？

提示：把环面还原成长方形，测地线在长方形里是直线；边的粘合意味着直线"走出一边就从对边回来"。

\section*{通往下一章的门}

这一章里，我们学会了在弯曲的世界里测量距离和方向。可是请注意一个被我们悄悄略过的前提：测量是精确的。绳子拉紧贴住球面，量出的弧长就是弧长；量角器读出的角度，就是角度。

现实中的测量从不这么乖。测量一座山的高度，今天测是 $8848.86$ 米，明天换一台仪器是 $8848.83$ 米；导航卫星报出你的位置，每次都有几米的飘忽；就连光速这样的常数，测量值也在小数点后面微微抖动。面对一堆"都对、又都不完全一样"的数字，哪一个是真相？还是说，"真相"这个词本身就需要重新定义？

从微分几何走向广义相对论，人类学会了描述弯曲但确定的空间；而下一章要面对的，是另一种更日常的弯曲——不确定性本身。让我们走进概率与统计的世界。
